%! Author = Kevin Lin
%! Date = 9/29/2026

% Preamble
\documentclass[11pt,a4paper,margin=1in]{article}

% Packages
\usepackage{amsmath}
\usepackage{amssymb}
\usepackage{enumerate}
\usepackage{pdfpages}

\title{HW 1}
\author{Kevin Lin}
\date{9/29/2026}

% Document
\begin{document}
\maketitle

\section{}
\begin{enumerate}[(a)]
    \item Given random variable $X$ where $a \leq X \leq b$, we can prove that 
    $X - E[X]$ is $\frac{(b-a)^2}{4}$-subgaussian. Let $Y = X - E[X]$ and the
    log MGF is $\psi(\lambda) = \log E[e^{\lambda Y}]$. Using the identity from 
    the hint, we can bound $\psi''(\lambda)$ as follows:
    \begin{gather*}
        \text{Let } M(\lambda) = E[e^{\lambda Y}] \\
        \psi(\lambda) = \log M(\lambda) \\
        \psi'(\lambda) = \frac{M'(\lambda)}{M(\lambda)} = \frac{E[Y e^{\lambda Y}]}{E[e^{\lambda Y}]} \\
        \psi''(\lambda) = \frac{M''(\lambda) M(\lambda) - (M'(\lambda))^2}{(M(\lambda))^2} = \frac{E[Y^2 e^{\lambda Y}] E[e^{\lambda Y}] - (E[Y e^{\lambda Y}])^2}{(E[e^{\lambda Y}])^2} \\
    \end{gather*}
    Since $Y = X - E[X]$ and subtracting a constant does not change variance:
    \[
        \psi''(\lambda) = \frac{E[Y^2 e^{\lambda Y}]}{E[e^{\lambda Y}]} -
        \left( \frac{E[Y e^{\lambda Y}]}{E[e^{\lambda Y}]} \right)^2 = Var(\tilde{X}) \leq \frac{(b-a)^2}{4}
    \]
    where $\tilde{X}$ is a random variable that satisfies $a \leq \tilde{X} \leq b$.
    Therefore, we have:
    \[
        \psi(\lambda) \leq \int_0^\lambda \int_0^t \frac{(b-a)^2}{4} ds dt = \frac{(b-a)^2 \lambda^2}{8}
    \]
    A mean zero random variable is $\sigma^2$-subgaussian if $\psi(\lambda) \leq \frac{\sigma^2 \lambda^2}{2}$. 
    Thus, we have:
    \begin{align*}
        \frac{\sigma^2 \lambda^2}{2} &= \frac{(b-a)^2 \lambda^2}{8} \\
        \sigma^2 &= \frac{(b-a)^2}{4}
    \end{align*}
    Thus, $X - E[X]$ is $\frac{(b-a)^2}{4}$-subgaussian.
    \item The constant $\frac{(b-a)^2}{4}$ is the best possible constant. To show
    this, we use the example of a random variable $X$ that takes on the values 
    $a$ and $b$ with equal probability. Then, we have:
    \begin{align*}
        E[X] &= \frac{a+b}{2} \\
        Var(X) &= E[X^2] - (E[X])^2 = \frac{a^2 + b^2}{2} - \left( \frac{a+b}{2} \right)^2 = \frac{(b-a)^2}{4}
    \end{align*}
    Let $Y = X - E[X]$. Then, we have:
    \begin{align*}
        E[Y] &= E[X - E[X]] = E[X] - E[X] = 0 \\
        Var(Y) &= Var(X) = \frac{(b-a)^2}{4}
    \end{align*}
    Now, suppose $Y$ is $\sigma^2$-subgaussian. Then
    \[
        \psi_Y(\lambda) = \log E[e^{\lambda Y}] \leq \frac{\sigma^2 \lambda^2}{2}
    \]
    Since $\psi_Y''(0) = Var(Y)$ and $\left( \frac{\sigma^2 \lambda^2}{2} \right)'' = \sigma^2$, 
    the subgaussian bound near $\lambda = 0$ implies 
    \[
        \frac{(b-a)^2}{4} \leq \sigma^2
    \]
    Thus, $Y = X - E[X]$ cannot be $\sigma^2$-subgaussian for any $\sigma^2 < \frac{(b-a)^2}{4}$. 
    Therefore, the constant $\frac{(b-a)^2}{4}$ cannot be improved in general.
\end{enumerate}

\section{}
Let $X$ be a random variable where $E[X] = 0$, $Var(X) \leq \sigma^2$, and $|X| \leq b$
almost surely for a finite constant $b$. For any $\lambda \in (-3/b, 3/b)$, we can
show that $\psi(\lambda) \leq \frac{\lambda^2\sigma^2}{2(1 - |\lambda| b / 3)}$
by using the Taylor exapnsion of $e^{\lambda X}$:
\begin{gather*}
    e^{\lambda X} = 1 + \lambda X + \frac{\lambda^2 X^2}{2!} + \frac{\lambda^3 X^3}{3!} + \cdots \\
    E[e^{\lambda X}] = 1 + \frac{\lambda^2 E[X^2]}{2!} + \frac{\lambda^3 E[X^3]}{3!} + \cdots \\
\end{gather*}
We can define an upper bound $\lambda^n X^n \leq |\lambda|^n |X|^n$, but since
$|X| \leq b$, we have $|X|^n \leq b^{n - 2}|X|^2$. So:
\begin{gather*}
    E[e^{\lambda X}] \leq 1 + E[X^2] \left( \frac{\lambda^2}{2!} + \frac{|\lambda|^3 b}{3!} + \frac{|\lambda|^4 b^2}{4!} + \cdots \right) \\
    E[X] = 0, \quad E[X^2] = Var(X) \leq \sigma^2 \therefore \\
    E[e^{\lambda X}] \leq 1 + \sigma^2 \left( \frac{\lambda^2}{2!} + \frac{|\lambda|^3 b}{3!} + \frac{|\lambda|^4 b^2}{4!} + \cdots \right) \\
    E[e^{\lambda X}] \leq 1 + \sigma^2 \left( \frac{\lambda^2}{2} \right) \left(1 + \frac{|\lambda| b}{3} + \frac{(|\lambda| b)^2}{12} + \cdots \right) \\
\end{gather*}
But we can recognize the series as a geometric series for the third term onward
since $\frac{|\lambda|^4b^2}{4!} \leq \frac{\lambda^2b^2}{2} \left( \frac{|\lambda|b}{3} \right)^2$.
Thus, we can write:
\[
    E[e^{\lambda X}] \leq 1 + \frac{\lambda^2\sigma^2}{2} \left(1 + \frac{|\lambda| b}{3} + \left(\frac{|\lambda| b}{3}\right)^2 + \cdots \right)
\]
We know that $|\lambda| < \frac{3}{b}$, so $\frac{|\lambda| b}{3} < 1$ and we can 
sum the geometric series using $r = \frac{|\lambda| b}{3}$:
\begin{gather*}
    E[e^{\lambda X}] \leq 1 + \frac{\lambda^2\sigma^2}{2} \left( \frac{1}{1 - |\lambda| b / 3} \right) \\
    E[e^{\lambda X}] \leq 1 + \frac{\lambda^2\sigma^2}{2(1 - |\lambda| b / 3)} \\
    log(1 + x) \leq x \therefore \\
    \log E[e^{\lambda X}] = \psi(\lambda) \leq \frac{\lambda^2\sigma^2}{2(1 - |\lambda| b / 3)}
\end{gather*}

\section{}
\begin{enumerate}[(a)]
    \item Since $X$ is a mean-zero random variable, and $X_+ = \max(X, 0)$, for any $t > 0$:
    \begin{gather*}
        P[X \geq t] = P[X_+ \geq t] \\
        P[X_+ \geq t] = P[X_+^p \geq t^p] \\
        \text{By Markov's inequality, } P[X_+^p \geq t^p] \leq \frac{E[X_+^p]}{t^p} \\
        P[X_+ \geq t] \leq \inf_{p > 0} \frac{E[X_+^p]}{t^p} \\
    \end{gather*}
    For the second inequality, using $p \rightarrow 0_+$, we get that for each fixed $x$:
    \begin{gather*}
        x_+^p \rightarrow \begin{cases}
            1, & x > 0 \\
            0, & x \leq 0
        \end{cases} \therefore \\
        X_+^p \rightarrow \mathbf{1}_{X > 0} \therefore \\
        \frac{E[X_+^p]}{t^p} \rightarrow \frac{P[X > 0]}{1} = P[X > 0] \\
    \end{gather*}
    Let $p = \lambda t$ for $\lambda > 0$. Then, since $t > 0$, we have $p > 0$, thus:
    \begin{gather*}
        \inf_{p > 0} \frac{E[X_+^p]}{t^p} \leq \frac{E[X_+^{\lambda t}]}{t^{\lambda t}} \\
        \frac{x_+^{\lambda t}}{t^{\lambda t}} = \left( \frac{x}{t} \right)^{\lambda t} = e^{\lambda t \log(x/t)} \\
        \log\left(\frac{x}{t}\right) \leq \frac{x}{t} - 1 \\
        \lambda t \log\left(\frac{x}{t}\right) \leq \lambda (x - t) \\
        \left(\frac{x}{t}\right)^{\lambda t} \leq e^{\lambda (x - t)} \therefore \\
        \frac{x_+^{\lambda t}}{t^{\lambda t}} \leq e^{\lambda (x - t)} \therefore \\
        \frac{E[X_+^{\lambda t}]}{t^{\lambda t}} \leq E[e^{\lambda (X - t)}] = e^{-\lambda t} E[e^{\lambda X}] \therefore \\
        \inf_{p > 0} \frac{E[X_+^p]}{t^p} \leq e^{-\lambda t} E[e^{\lambda X}] \quad \forall \lambda > 0
    \end{gather*}
    For $\lambda = 0$, we have from the hint $e^{-\lambda t} E[e^{\lambda X}] = 1 \geq P[X > 0]$. 
    Since the inequality holds over all $\lambda \geq 0$, we can take the infimum over $\lambda \geq 0$:
    \[
        \inf_{p > 0} \frac{E[X_+^p]}{t^p} \leq \inf_{\lambda \geq 0} e^{-\lambda t} E[e^{\lambda X}]
    \]
    which gives us the full inequality:
    \[
        P[X \geq t] \leq \inf_{p > 0} \frac{E[X_+^p]}{t^p} \leq \inf_{\lambda \geq 0} e^{-\lambda t} E[e^{\lambda X}]
    \]
    \item Let $X = \epsilon Y$ where $\epsilon$ is a Rademacher random variable 
    independent of $Y$, where $P[Y \geq y] = y^{-\alpha}$ for $y \geq 1$ and $\alpha > 1$.
    For $t \geq e$, the exact tail probability $P[X \geq t]$ is:
    \begin{gather*}
        P[X \geq t] = P[\epsilon Y \geq t] = P[\epsilon = 1, Y \geq t] + P[\epsilon = -1, Y \geq t] \\
        P[X \geq t] = \frac{1}{2} P[Y \geq t] + 0 = \frac{1}{2} t^{-\alpha} \\
    \end{gather*}
    For the first infimum, the exact tail probability $P[X \geq t]$ is:
    \begin{gather*}
        X_+ = \begin{cases}
            Y, & \epsilon = 1 \\
            0, & \epsilon = -1
        \end{cases} \therefore \\
        E[X_+^p] = \frac{1}{2} E[Y^p] + \frac{1}{2} E[0^p] = \frac{1}{2} E[Y^p] \\
        E[Y^p] = \int_1^\infty y^p \alpha y^{-\alpha - 1} dy = \frac{\alpha}{\alpha - p} \quad \text{for } 0 < p < \alpha \therefore \\
        \inf_{p > 0}\frac{E[X_+^p]}{t^p} = \frac{\alpha}{2(\alpha - p)t^p} \\
    \end{gather*}
    We need to minimize this:
    \begin{gather*}
        \log\left(\frac{\alpha}{2(\alpha - p)t^p}\right) = \log\left(\frac{\alpha}{2}\right) - \log(\alpha - p) - p \log(t) \\
        \frac{d}{dp} \left( \log\left(\frac{\alpha}{2}\right) - \log(\alpha - p) - p \log(t) \right) = 0 \\
        \frac{1}{\alpha - p} - \log(t) = 0 \therefore \\
        p = \alpha - \frac{1}{\log(t)} \therefore \\
        \inf_{p > 0}\frac{E[X_+^p]}{t^p} = \frac{\alpha}{2(1/\log(t))t^{\alpha - 1/\log(t)}} = \frac{\alpha \log(t)}{2} t^{-\alpha + 1/\log(t)} = \frac{e\alpha\log(t)}{2et^{\alpha}}\\
    \end{gather*}
    Compared to the original tail probability $P[X \geq t] = \frac{1}{2} t^{-\alpha}$,
    the overestimate factor is:
    \[
        \frac{e\alpha\log(t)}{2et^{\alpha}} \cdot \frac{2t^{\alpha}}{1} = e\alpha\log(t)
    \]
    For the second infimum, we have:
    \begin{gather*}
        E[e^{\lambda X}] = \frac{1}{2} E[e^{\lambda Y}] + \frac{1}{2} E[e^{-\lambda Y}] \\
    \end{gather*}
    But $E[e^{\lambda Y}] = \int_1^\infty e^{\lambda y} \alpha y^{-\alpha - 1} dy$ 
    diverges for any $\lambda > 0$, so we can only consider $\lambda = 0$. Thus, the second infimum is:
    \[
        \inf_{\lambda \geq 0} e^{-\lambda t} E[e^{\lambda X}] = 1
    \]
    Compared to the original tail probability $P[X \geq t] = \frac{1}{2} t^{-\alpha}$,
    the overestimate factor is:
    \[
        1 \cdot \frac{2t^{\alpha}}{1} = 2t^{\alpha}
    \]
\end{enumerate}

\section{}
\begin{enumerate}[(a)]
    \item If $X = 0$, then for $E[e^{X^2/K^2}] = E[e^{0}] = 1 \leq 2$, therefore
    $||X||_{\psi_2} = \inf\{K > 0 : E[e^{X^2/K^2}] \leq 2\} = 0$. Suppose $X \neq 0$,
    then for $K \rightarrow 0^+$, $e^{X^2/K^2} \rightarrow \infty$ and consequently
    $E[e^{X^2/K^2}] \rightarrow \infty > 2$. Thus, $||X||_{\psi_2} > 0$, which
    contradicts $||X||_{\psi_2} = 0$. Therefore, $||X||_{\psi_2} = 0$ if and only if $X = 0$.
    \item $||cX||_{\psi_2} = |c|||X||_{\psi_2}$. If $c = 0$, both sides are zero.
    If $c \neq 0$, then:
    \begin{gather*}
        ||cX||_{\psi_2} = \inf\{K > 0 : E[e^{(cX)^2/K^2}] \leq 2\} \\
        ||cX||_{\psi_2} = \inf\{K > 0 : E[e^{X^2/(K/|c|)^2}] \leq 2\}, \quad K' = K/|c| \\
        ||cX||_{\psi_2} = |c| \inf\{K' > 0 : E[e^{X^2/(K')^2}] \leq 2\} = |c|||X||_{\psi_2}
    \end{gather*}
    \item Let $a > ||X||_{\psi_2}$ and $b > ||Y||_{\psi_2}$. Then, we have:
    \begin{gather*}
        E[e^{X^2/a^2}] \leq 2, \quad E[e^{Y^2/b^2}] \leq 2 \\
        E[e^{(X+Y)^2/(a+b)^2}] = E[e^{((a/(a+b))(X/a) + (b/(a+b))(Y/b))^2}] \\
        \text{Let } \theta = \frac{a}{a+b}, \quad 1 - \theta = \frac{b}{a+b} \\
        \left(\theta(X/a) + (1 - \theta)(Y/b)\right)^2 \leq (\theta)(X/a)^2 + (1 - \theta)(Y/b)^2 \\
        E[e^{(X+Y)^2/(a+b)^2}] \leq E[e^{(\theta)(X/a)^2 + (1 - \theta)(Y/b)^2}] \\
        E[e^{(X+Y)^2/(a+b)^2}] \leq E[e^{(\theta))(X/a)^2} e^{(1 - \theta)(Y/b)^2}] \\
        E[e^{(X+Y)^2/(a+b)^2}] \leq E[e^{(\theta)(X/a)^2}] E[e^{(1 - \theta)(Y/b)^2}] \\
        E[e^{(X+Y)^2/(a+b)^2}] \leq (E[e^{(X/a)^2}])^{\theta} (E[e^{(Y/b)^2}])^{1 - \theta} \\
        E[e^{(X+Y)^2/(a+b)^2}] \leq (E[e^{(X/a)^2}]E[e^{(Y/b)^2}])^{\theta + (1 - \theta))} \\
        E[e^{(X+Y)^2/(a+b)^2}] \leq 2^1 = 2, \therefore\\
        ||X+Y||_{\psi_2} \leq a + b = ||X||_{\psi_2} + ||Y||_{\psi_2}
    \end{gather*}
\end{enumerate}

\section{}
\begin{enumerate}[(a)]
    \item Let $X$ be a random variable where $E[X] = 0$, $Var[X] = \sigma^2$, and
    $X \leq b$ with probability 1. Then, we can show that for all $\lambda \geq 0$:
    \[
        \log E[e^{\lambda X}] \leq \frac{\sigma^2(e^{\lambda b} - 1 - \lambda b)}{b^2}
    \]
    by starting with $e^{\lambda X} = 1 + \lambda X + (e^{\lambda X} - 1 - \lambda X)$
    and bounding the last term. Because $X \leq b$ and $\lambda \geq 0$, $\lambda X \leq \lambda b$.
    From the hint, we can then know that:
    \begin{gather*}
        \frac{e^{\lambda X} - 1 - \lambda X}{(\lambda X)^2} \leq \frac{e^{\lambda b} - 1 - \lambda b}{(\lambda b)^2} \\
        e^{\lambda X} - 1 - \lambda X \leq \frac{X^2(e^{\lambda b} - 1 - \lambda b)}{b^2} \\
        e^{\lambda X} \leq 1 + \lambda X + \frac{X^2(e^{\lambda b} - 1 - \lambda b)}{b^2} \\
        E[e^{\lambda X}] \leq 1 + \lambda E[X] + \frac{E[X^2](e^{\lambda b} - 1 - \lambda b)}{b^2} \\
        E[X] = 0, \quad E[X^2] = Var[X] = \sigma^2 \therefore \\
        E[e^{\lambda X}] \leq 1 + \frac{\sigma^2(e^{\lambda b} - 1 - \lambda b)}{b^2} \\
        \log(E[e^{\lambda X}]) \leq \log\left(1 + \frac{\sigma^2(e^{\lambda b} - 1 - \lambda b)}{b^2}\right) \\
        \log\left(1 + \frac{\sigma^2(e^{\lambda b} - 1 - \lambda b)}{b^2}\right) \leq \frac{\sigma^2(e^{\lambda b} - 1 - \lambda b)}{b^2} \therefore \\
        \log(E[e^{\lambda X}]) \leq \frac{\sigma^2(e^{\lambda b} - 1 - \lambda b)}{b^2}
    \end{gather*}
    \item Now, with $X_1, \dots, X_n$ independent random variables equal in law
    to $X$, for all $t \geq 0$, we can show that:
    \[
        P\left[\sum_{i=1}^n X_i \geq t\right] \leq \exp\left(-\frac{n\sigma^2}{b^2} h\left(\frac{bt}{n\sigma^2}\right)\right)
    \]
    where $h(x) = (1+x)\log(1+x) - x$ by using the Chernoff bound. From part (a)
    and because each $X_i$ is independent and equal in law to $X$, we have:
    \begin{gather*}
        E[e^{\lambda \sum_{i=1}^n X_i}] = E\left[\prod_{i=1}^n e^{\lambda X_i}\right] = (E[e^{\lambda X}])^n \\
        \log E[e^{\lambda \sum_{i=1}^n X_i}] = n \log E[e^{\lambda X}] \leq \frac{n\sigma^2(e^{\lambda b} - 1 - \lambda b)}{b^2} \therefore \\
        E[e^{\lambda \sum_{i=1}^n X_i}] \leq \exp\left(\frac{n\sigma^2(e^{\lambda b} - 1 - \lambda b)}{b^2}\right)
    \end{gather*}
    Now using the Chernoff bound, we have:
    \[
        P[\sum_{i=1}^n X_i \geq t] \leq \exp\left(-\lambda t + \frac{n\sigma^2(e^{\lambda b} - 1 - \lambda b)}{b^2}\right) \\
    \]
    Minimizing the right-hand side over $\lambda \geq 0$, we have:
    \begin{gather*}
        \frac{d}{d\lambda} \left(-\lambda t + \frac{n\sigma^2(e^{\lambda b} - 1 - \lambda b)}{b^2}\right) = 0 \\
        -t + \frac{n\sigma^2(e^{\lambda b}b - b)}{b^2} = 0 \therefore \\
        e^{\lambda b} = 1 + \frac{bt}{n\sigma^2} \therefore \\
        \lambda^* = \frac{1}{b}\log\left(1 + \frac{bt}{n\sigma^2}\right)
    \end{gather*}
    Let $x = \frac{bt}{n\sigma^2}$. Then, we have:
    \begin{gather*}
        \lambda^* = \log(1 + x), \quad e^{\lambda^* b} = 1 + x, \quad t = \frac{n\sigma^2 x}{b} \\
        \therefore \\
        -\lambda^* t = -\frac{n\sigma^2 x \log(1+x)}{b^2} \\
        \frac{n\sigma^2(e^{\lambda b} - 1 - \lambda b)}{b^2} = \frac{n\sigma^2((1+x) - 1 - \log(1+x))}{b^2} = \frac{n\sigma^2(x - \log(1+x))}{b^2} \\
        \therefore
    \end{gather*}
    \begin{align*}
        -\lambda t + \frac{n\sigma^2(e^{\lambda b} - 1 - \lambda b)}{b^2} &= -\frac{n\sigma^2 x \log(1+x)}{b^2} + \frac{n\sigma^2(x - \log(1+x))}{b^2} \\
        -\lambda t + \frac{n\sigma^2(e^{\lambda b} - 1 - \lambda b)}{b^2} &= -\frac{n\sigma^2}{b^2}(x\log(1 + x) + x - \log(1 + x)) \\
        -\lambda t + \frac{n\sigma^2(e^{\lambda b} - 1 - \lambda b)}{b^2} &= -\frac{n\sigma^2}{b^2}((1+x)\log(1+x) - x) \\
        -\lambda t + \frac{n\sigma^2(e^{\lambda b} - 1 - \lambda b)}{b^2} &= -\frac{n\sigma^2}{b^2} h(x) \\
    \end{align*}
    Then finally, substituting back $x = \frac{bt}{n\sigma^2}$, we have:
    \[
        P\left[\sum_{i=1}^n X_i \geq t\right] \leq \exp\left(-\frac{n\sigma^2}{b^2} h\left(\frac{bt}{n\sigma^2}\right)\right)
    \]
\end{enumerate}

\section{}
\begin{enumerate}[(a)]
    \item We can show that if $X_1, \dots, X_n$ are mean-zero and each $X_i$ is
    $\sigma_i^2$-subgaussian, then without any assumption on their dependence, \\
    $\sum_{i=1}^n X_i$ is $\sum_{i=1}^n \sigma_i^2$-subgaussian.
    Let $A = \sum_{i=1}^n \sigma_i$ and weights $w_i = \frac{\sigma_i}{A}$. 
    Then, we have:
    \begin{gather*}
        \sum_{i=1}^n w_i = 1, \quad w_i \geq 0 \\
        w_i\frac{X_i}{\sigma_i} = \frac{X_i}{A} \therefore \\
        \sum_{i=1}^n w_i\frac{X_i}{\sigma_i} = \sum_{i=1}^n \frac{X_i}{A} \\
        \sum_{i=1}^n X_i = A \sum_{i=1}^n w_i\frac{X_i}{\sigma_i} \therefore
    \end{gather*}
    Because $w_i \geq 0$ and $\sum_{i=1}^n w_i = 1$, we can use Jensen's inequality
    and the convexity of $e^x$ to show that:
    \begin{gather*}
        e^{\lambda A \sum_{i=1}^n w_i(X_i/\sigma_i)} \leq \sum_{i=1}^n w_i e^{\lambda A (X_i/\sigma_i)} \\
        E[e^{\lambda A \sum_{i=1}^n X_i}] \leq E\left[\sum_{i=1}^n w_i e^{\lambda A (X_i/\sigma_i)}\right] = \sum_{i=1}^n w_i E[e^{\lambda A (X_i/\sigma_i)}] \\
    \end{gather*}
    Since $X_i$ is $\sigma_i^2$-subgaussian, we have:
    \begin{gather*}
        E[e^{tX_i}] \leq e^{t^2\sigma_i^2/2}, \text{ if } t = \lambda A / \sigma_i \\
        E[e^{\lambda A (X_i/\sigma_i)}] \leq e^{(\lambda A / \sigma_i)^2\sigma_i^2/2} = e^{\lambda^2 A^2/2} \\
        E[e^{\lambda A \sum_{i=1}^n X_i}] \leq \sum_{i=1}^n w_i e^{\lambda^2 A^2/2} = e^{\lambda^2 A^2/2} \\
    \end{gather*}
    Recall $A = \sum_{i=1}^n \sigma_i$, so we have:
    \[
        E[e^{\lambda \sum_{i=1}^n X_i}] \leq \exp\left(\frac{\lambda^2}{2} \left(\sum_{i=1}^n \sigma_i\right)^2\right)  
    \]
    thus $\sum_{i=1}^n X_i$ is $\left(\sum_{i=1}^n \sigma_i\right)^2$-subgaussian.
    This bound can not be improved in general. Let $\epsilon$ be a Rademacher
    random variable and let $X_i = \sigma_i \epsilon$. Then, we have:
    \begin{gather*}
        \sum_{i=1}^n X_i = \sum_{i=1}^n \sigma_i \epsilon = \left(\sum_{i=1}^n \sigma_i\right) \epsilon \therefore \\
        Var(S) = \left(\sum_{i=1}^n \sigma_i\right)^2
    \end{gather*}
    Any mean-zero random variable is $\sigma^2$-subgaussian if $\sigma^2 \geq Var(X)$, so
    $\sum_{i=1}^n X_i$ is $\left(\sum_{i=1}^n \sigma_i\right)^2$-subgaussian, and the bound
    cannot be improved in general.
    \item Let $k \geq 2$ and $\epsilon_1, \dots, \epsilon_k$ be independent Rademacher 
    random variables. For every nonempty subset $S \subseteq \{1, \dots, k\}$, let 
    $X_S = \prod_{i \in S} \epsilon_i$, giving $n = 2^k - 1$ random variables and
    their respective sum $W = \sum_S X_S$.
    \begin{enumerate}[i.]
        \item We can show each $X_S$ is a Rademacher random variable. Picking any
        $j \in S$, $X_S = \epsilon_j \prod_{i \in S, i \neq j} \epsilon_i$. Since 
        $\epsilon_j$ is a Rademacher variable, by conditioning on the other $\epsilon_i$'s, 
        $\prod_{i \in S, i \neq j} \epsilon_i$ is a constant, and thus $X_S$ is 
        a Rademacher variable. $X_S$ and $X_T$ are independent for any $S \neq T$:
        \begin{gather*}
            X_SX_T = \left(\prod_{i \in S} \epsilon_i\right)\left(\prod_{j \in T} \epsilon_j\right) \\
            X_SX_T = \prod_{i \in S \triangle T} \epsilon_i
        \end{gather*}
        Because $S \neq T$, $S \triangle T$ is nonempty, and thus $X_SX_T$ is a Rademacher
        variable with expectation 0 from the prior proof. Thus, $Cov(X_S, X_T) = E[X_SX_T] - E[X_S]E[X_T] = 0 - 0 = 0$, 
        and $X_S$ and $X_T$ are uncorrelated, proving their independence.
        \item Recognize:
        \begin{gather*}
            \prod_{i=1}^k (1 + \epsilon_i) = 1 + \sum_{S \subseteq \{1, \dots, k\}} X_S = 1 + W \therefore \\
            W = \prod_{i=1}^k (1 + \epsilon_i) - 1
        \end{gather*}
        Because each $\epsilon_i$ is either $-1$ or 1, a single $\epsilon_i = -1$ 
        will make the product 0, making $W = -1$. Only if all $\epsilon_i = 1$ 
        will the product be nonzero, making $W = 2^k - 1 = n$. Since all $\epsilon_i$ 
        are independent, $P(\epsilon_1 = 1, \dots, \epsilon_k = 1) = P(W = n) = (1/2)^k$,
        thus $P(W = n) = 1/(n + 1)$. $n$ is also the largest possible value of $W$,
        so $P(W \geq n) = P(W = n) = 1/(n + 1)$.
        \item If $W$ were $\sigma^2$-subgaussian, then we would have:
        \begin{gather*}
            P(W \geq t) \leq e^{-t^2/(2\sigma^2)} \\
            P(W \geq n) \leq e^{-n^2/(2\sigma^2)} \therefore \\
            \frac{1}{n + 1} \leq e^{-n^2/(2\sigma^2)} \\
            -\log(n + 1) \leq -\frac{n^2}{2\sigma^2} \\
            \log(n + 1) \geq \frac{n^2}{2\sigma^2} \\
            \sigma^2 \geq \frac{n^2}{2\log(n + 1)}
        \end{gather*}
        Thus, $W$ cannot be $\sigma^2$-subgaussian for any $\sigma^2 < \frac{n^2}{2\log(n + 1)}$.
        Chebyshev's inequality gives us:
        \[
            P(W \geq n) \leq P(|W| \geq n) \leq \frac{Var(W)}{n^2} = \frac{2^k - 1}{n^2} = \frac{n}{n^2} = \frac{1}{n} \\
        \]
        overestimating the tail probability by a factor of $\frac{n + 1}{n}$.
        Hoeffding's inequality gives us:
        \[
            P(W \geq n) \leq e^{-2n^2/(4n)} = e^{-n/2} \\
        \]
        which overestimates the tail probability by a factor of $(n + 1)e^{-n/2}$.
        But this is less than 1 for any $n \geq 3$, so it is not a valid upper
        bound. Pairwise independence is not enough to justify the Hoeffding bound,
        but gives the correct variance behavior (just not the full subgaussian tail).
    \end{enumerate}
\end{enumerate}

\noindent 
\textbf{AI Usage:} All homework problems were written in VSCode and compiled using LaTeX, with 
assistance from ChatGPT and Github Copilot inline suggestions.

\end{document}